\section*{الفصل \ref{ch:b2:genome-diversification} --- الطفرات وتنويع المورثوم}

\begin{solution}{exo:b2:genome-diversification:1}
GAG: صامتة (ولا تزال غلوتامات)، وهي انتقال (A$\to$G). وGTA: خاطئة
المعنى (فالين)، وهي تحويل (A$\to$T). وTAA: عديمة المعنى (إيقاف)،
وهي تحويل (G$\to$T). وإقحام أساس: إزاحة إطار، تعيد كتابة كل رامزة
بعدها.
\end{solution}

\begin{solution}{exo:b2:genome-diversification:2}
$U = 5\times 10^{-10}\times 4.6\times 10^{6} = 2.3\times 10^{-3}$: أي
إن انقسامًا من كل 430 تقريبًا ينتج \omterm{def:b2:genome-diversification:mutation}{طفرة}. وفي $10^{9}$ خلية:
$2.3\times 10^{6}$ \omterm{def:b2:genome-diversification:mutation}{طفرة} جديدة في الجيل.
\end{solution}

\begin{solution}{exo:b2:genome-diversification:3}
التحول: التقاط DNA حر من خلايا متحللة، ويلزمه متلقٍّ كفؤ؛ وينقل
قطعًا صبغية (مورّثات المحفظة عند المكورات الرئوية). والاقتران: تماس
بين خلية وأخرى عبر هلب، ويلزمه \omterm{def:b2:unicellular-diversity:prokaryote}{بلازميد} اقتراني في الواهب؛ وينقل
\omterm{def:b2:unicellular-diversity:prokaryote}{البلازميدات} (مورّثات المقاومة)، وينقل من سلالة Hfr مورّثات صبغية.
\omterm{def:b2:genome-diversification:hgt}{والنقل بالعاثيات}: عاثية حزمت DNA مضيفًا؛ وتنقل قطعة عشوائية (نقل
عام) أو المورّثات المجاورة لموضع العاثية المدمجة (نقل خاص).
\end{solution}

\begin{solution}{exo:b2:genome-diversification:4}
متعدد الصيغة الذاتي: طقوم زائدة من نوع واحد، بعروس ثنائية الصيغة
(البطاطا رباعية الصيغة، وكثير من النفل الرباعي الذاتي). ومتعدد
الصيغة المغاير: طقما نوعين مجموعان بالتهجين ثم مضاعفان (قمح الخبز،
والقطن، والتبغ). والهجين الثنائي الصيغة AB ليس فيه صبغيان
متشابهان، فلا يستطيع أي منهما التزاوج في التقسم المنصّف الأول؛
فتنفصل الصبغيات عشوائيًا وتكون الأعراس غير متوازنة وغير قابلة للحياة.
\end{solution}

\begin{solution}{exo:b2:genome-diversification:5}
$2\times 3.2\times 10^{9}\times 1.2\times 10^{-8} = 77$ \omterm{def:b2:genome-diversification:mutation}{طفرة} جديدة.
وفي التسلسل المرمّز: $77\times 0.015 = 1.2$؛ ومنها ما يغير حمضًا
أمينيًا: $1.2\times 0.7 = 0.8$ --- أي إن معظم الأطفال يحملون نحو
استبدال جديد واحد في الحموض الأمينية.
\end{solution}

\begin{solution}{exo:b2:genome-diversification:6}
الأحداث: $m(N - N_0) = 2\times 10^{-8}\times 10^{9} = 20$. والخلايا
الطافرة: $mN\ln(N/N_0) = 20\times \ln 10^{6} = 20\times 13.8 = 276$.
و$p_0 = e^{-20} \approx 2\times 10^{-9}$: أي إن كل مزرعة فيها
طافرات، فلا يقيس عدّ المزارع الخالية شيئًا. ويلزم أن يكون $m(N -
N_0) \approx 1$، أي $N \approx 5\times 10^{7}$، حيث $p_0 \approx
0.37$ وتعطي بضع عشرات من المزارع $m$ في حدود عامل قدره اثنان.
\end{solution}

\begin{solution}{exo:b2:genome-diversification:7}
لا مكان لليوراسيل في DNA، فتستطيع غليكوزيلاز اليوراسيل--DNA إزالة
كل يوراسيل تجده دون لبس. أما الثيمين فأساس عادي: فبعد \omterm{prop:b2:genome-diversification:sources}{نزع الأمين} من
5-ميثيل سيتوزين يُرى عدم التوافق G$\cdot$T، لكن جهاز الإصلاح لا
سبيل له إلى معرفة أي السلسلتين خطأ، فيصحح G في نصف الحالات، فيثبّت
انتقال C$\to$T. ولذلك تطفر مواقع CpG المُمثيلة أسرع بنحو عشر مرات
من سواها، وقد تحولت على مدى الزمن التطوري إلى TpG وCpA: فمورثومات
الثدييات لا تحوي إلا نحو خمس ما تتنبأ به تركيبة الأسس من CpG.
\end{solution}

\begin{solution}{exo:b2:genome-diversification:8}
\qty{46}{kb} في الدقيقة. والمواضع: 368 و782 و1150
و\qty{1840}{kb}؛ والمسافات 414 و368 و\qty{690}{kb}.
\end{solution}

\begin{solution}{exo:b2:genome-diversification:9}
تتزاوج النسخة الثانية من أحد الصبغيين مع النسخة الأولى من الآخر؛
ويعطي تعابر بينهما كروماتيدة تحمل النسخة~1 والنسخة~2 والنسخة~2$'$
(ثلاث مورّثات) وكروماتيدة لا تحمل إلا النسخة~1$'$. والشخص الذي ورث
صبغيًا ذا مورّثة واحدة من كل من أبويه له مورّثتا $\alpha$ بدل أربع
($-\alpha/-\alpha$): فتُصنع سلاسل $\alpha$ بنصف المعدل، وتكون
الكريات الحمر صغيرة شاحبة، مع فقر دم خفيف أو بلا فقر دم --- وهي
سمة الثلاسيميا $\alpha$.
\end{solution}

\begin{solution}{exo:b2:genome-diversification:10}
الأعراس: $n = 21$ (طقم واحد من كل من A وB وD). والهجن: AB وفيه 14
صبغيًا؛ وABD وفيه 21. وقمح الخبز $\times$ قمح الإمر: AABBD وفيه 35
صبغيًا --- إذ يتزاوج طقما A وB، ولا شريك لطقم D، فيكون النبات عقيمًا
في جله. وكل مضاعفة أعطت كل صبغي متماثلًا مطابقًا، فتتكون الثنائيات
وتكون الأعراس متوازنة؛ أما التهجين الرجعي مع أحد الأبوين فينتج
نباتات فيها طقم واحد غير متزاوج، وأعراسها غير متوازنة، فيكون سداسي
الصيغة معزولًا تكاثريًا عن أسلافه --- أي نوعًا في خطوة واحدة.
\end{solution}

\begin{solution}{exo:b2:genome-diversification:11}
نمو لوجستي: $R(t) = N/\bigl(1 + (N - 1)e^{-\gamma Nt}\bigr)$.
ويبلغ نصف الجمهرة حين $(N - 1)e^{-\gamma Nt} = 1$: أي $t = \ln(N -
1)/\gamma N = 20.7/0.5 = \qty{41}{h}$. ومن دون الصادّ تنمو الخلايا
الخالية من \omterm{def:b2:unicellular-diversity:prokaryote}{البلازميد} أسرع بنسبة \qty{5}{\%}، فتتناقص نسبة الحاملات
إلى غير الحاملات على صورة $e^{-st}$ مع $s = 0.05\ln 2 =
\qty{0.035}{h^{-1}}$ في الساعة (\omterm{def:b2:unicellular-diversity:division}{وزمن الجيل} \qty{1}{h}): فالانتقال من
\qty{99}{\%} إلى \qty{1}{\%} يستغرق $\ln(99^2)/0.035 = \qty{260}{h}$،
أي أحد عشر يومًا --- والنقل المستمر يستطيع إبقاء \omterm{def:b2:unicellular-diversity:prokaryote}{البلازميد} إلى ما لا
نهاية عند تواتر منخفض، مهيّأً للجرعة التالية من الدواء.
\end{solution}

\begin{solution}{exo:b2:genome-diversification:12}
بيّن \omterm{thm:b2:genome-diversification:luria}{اختبار التذبذب} أن طفرات المقاومة تنشأ قبل العامل المنتخِب
ومن دونه، في أزمنة عشوائية، \omterm{def:b2:genome-diversification:mutation}{فالطفرة} عمياء عما قد ينفع؛ وبيّن الزرع
بالنسخ الشيء نفسه دون تعريض الخلايا المنتخبة البتة. أما التطور فليس
عشوائيًا لأن الانتخاب يفرز هذه المتغيرات العمياء بعواقبها، باطراد
وتراكم. والموضع الوحيد الذي تُشكَّل فيه العشوائية نفسها هو المعدل:
فالسلالات ذات الأخطاء الكثيرة أو القليلة جدًا تخسر، \omterm{thm:b2:genome-diversification:rate}{فمعدل الطفر} حل
وسط متكيّف --- لكن ما تفعله كل \omterm{def:b2:genome-diversification:mutation}{طفرة} يبقى غير متوقع وغير مختار.
\end{solution}

\begin{solution}{pb:b2:genome-diversification:1}
\textbf{1.} المجموع 119، والمتوسط $5.95$؛ ومتوسط المربعات
$11489/20 = 574$؛ والتباين $574 - 35 = 540$.
\textbf{2.} المجموع 324، والمتوسط $16.2$؛ ومتوسط المربعات
$5462/20 = 273$؛ والتباين $273 - 262 = 11$.
\textbf{3.} عينات المزرعة الواحدة بواسونية (فتباينها قريب من
متوسطها): فهي سحوب عشوائية من كسر طافر واحد ثابت. أما المزارع
المستقلة فتباينها تسعون ضعف متوسطها: فطافراتها نشأت في أزمنة مختلفة
عشوائية خلال النمو، وكلما كانت \omterm{def:b2:genome-diversification:mutation}{الطفرة} أبكر كانت نسيلتها أكبر.
\textbf{4.} جائزة كبرى: \omterm{def:b2:genome-diversification:mutation}{طفرة} في أحد الانقسامات الأولى، نمت نسيلتها
بعد ذلك مع المزرعة حتى بلغت مئة خلية.
\textbf{5.} أربع عشرة من عشرين: $p_0 = 0.70$؛ ومنه $m(N - N_0) = -\ln 0.70
= 0.36$؛ و$m = 0.36/2\times 10^{8} = \qty{1.8e-9}{}$ في الانقسام.
\textbf{6.} $1.8\times 10^{-9}\times 2\times 10^{8}\times \ln(2\times
10^{6}) = 0.36\times 14.5 = 5.2$ خلية طافرة؛ والمرصود $5.95$ ---
وهو توافق ضمن ضجيج عشرين مزرعة.
\textbf{7.} لم تخلُ أي عينة من الطافرات (فكلها تتشارك طافرات المزرعة
الأم)، فالمقدار $p_0 = 0$ لا يعطي تقديرًا؛ والعينات تقيس الكسر الطافر
لمزرعة واحدة، وهو يتوقف على متى وقعت طفراتها مصادفةً.
\textbf{8.} $u = m/20 = \qty{9e-11}{}$ لكل زوج أسس في النسخة الواحدة.
\textbf{9.} $U = 9\times 10^{-11}\times 4.6\times 10^{6} =
\qty{4.1e-4}{}$.
\textbf{10.} $10^{12}$ خلية $\times$ $4.1\times 10^{-4}$ $= 4.1\times
10^{8}$ \omterm{def:b2:genome-diversification:mutation}{طفرة} جديدة.
\textbf{11.} $3\times 4.6\times 10^{6} = 1.4\times 10^{7}$ تغير
ممكن؛ وينشأ كل منها $4.1\times 10^{8}/1.4\times 10^{7} \approx 30$
مرة.
\textbf{12.} كل تغير في أساس مفرد، بما فيه كل تغير يمنح المقاومة
لدواء تهزمه \omterm{def:b2:genome-diversification:mutation}{طفرة} واحدة، حاضر في عشرات الخلايا قبل إضافة الدواء؛
والدواء لا يفعل إلا أن ينتخبها.
\textbf{13.} تقع طفرتان مستقلتان في الخلية نفسها عند
$m^2 \approx (2\times 10^{-9})^2 = 4\times 10^{-18}$ في الانقسام: أي
في $10^{12}$ خلية، $4\times 10^{-6}$ في الجيل --- فالخلية المقاومة
مقاومة مضاعفة لا توجد سلفًا عمليًا، ولهذا يُعالج السل بعدة أدوية
دفعة واحدة.
\textbf{14.} بوضع $a = N - 1$ و$E = e^{-\gamma Nt}$: يكون $R = N/(1 +
aE)$، و$\mathrm{d}R/\mathrm{d}t = \gamma N\cdot NaE/(1 + aE)^2$؛
و$\gamma R(N - R) = \gamma\,\frac{N}{1 + aE}\cdot\frac{NaE}{1 + aE}$،
وهو نفسه؛ و$R(0) = N/(1 + N - 1) = 1$.
\textbf{15.} $(N - 1)E = 1$: ومنه $t = \ln(N - 1)/\gamma N = 20.7/0.5 =
\qty{41}{h}$.
\textbf{16.} $1 + aE = 1/0.99$، و$aE = 0.0101$: ومنه $t = \ln(10^{9}/0.0101)
/0.5 = (20.7 + 4.6)/0.5 = \qty{51}{h}$.
\textbf{17.} $\gamma (N/2)^2 = \gamma N\cdot N/4 = 0.5\times 2.5\times
10^{8} = 1.25\times 10^{8}$ نقلة في الساعة.
\textbf{18.} معدلا النمو $\ln 2 = \qty{0.693}{h^{-1}}$
و\qty{0.762}{h^{-1}}؛ وتنمو نسبة الطافرة إلى الحساسة على صورة $e^{st}$
مع $s = \qty{0.069}{h^{-1}}$؛ والانتقال من $1/N$ إلى 1 يستغرق $\ln N/s =
20.7/0.069 = \qty{300}{h}$، أي اثني عشر يومًا.
\textbf{19.} النقل يتناسب مع جداء الحاملات في غير الحاملات ويجري
بمعدل اللقاءات لا بمعدل الانقسامات؛ أما التناسل فيحده تفوق الطافرة
في النمو، وهو ضئيل. \omterm{def:b2:unicellular-diversity:prokaryote}{فالبلازميد} يعبر الجمهرة في يومين، والطافرة
المفضّلة في أسبوعين.
\textbf{20.} $U = 4.1\times 10^{-2}$؛ والضارة: $0.041\times 0.88\times
0.3 = 0.011$ في الجيل.
\textbf{21.} المتوسط $2.2$: ومنه $e^{-2.2} = 0.11$، أي خلف واحد من
كل تسعة يسلم.
\textbf{22.} النمط البري: $1.1\times 10^{-4}$ في الجيل، أي $0.022$ في
200 جيل: ومنه $e^{-0.022} = 0.98$.
\textbf{23.} المتوسط $520$ خلية طافرة؛ و$p_0 = e^{-36} \approx
10^{-16}$: أي لا مزرعة خالية.
\textbf{24.} تحت انتخاب شديد متغير (صادات مستشفى متناوبة) يرجح
العرض المئوي من المتغيرات الجديدة على العبء، وتستقل السلالات
المطفِّرة المقاومات التي تولدها. أما في محيط مستقر فلا شيء يُكسب
وثمة \omterm{def:b2:genome-diversification:mutation}{طفرة} ضارة كل مئة جيل تُخسر، فتُنقّى سلالات المطفِّرات.
\textbf{25.} $m = \qty{1.8e-9}{}$ في الانقسام نحو المقاومة، و$u =
\qty{9e-11}{}$ لكل زوج أسس في النسخة الواحدة، و$U = \qty{4.1e-4}{}$
\omterm{def:b2:genome-diversification:mutation}{طفرة} لكل مورثوم في النسخة الواحدة؛ ويبلغ \omterm{def:b2:unicellular-diversity:prokaryote}{البلازميد} نصف الأمعاء في
\qty{41}{h} و\qty{99}{\%} منها في \qty{51}{h}.
\end{solution}
